\documentclass[../main.tex]{subfiles}
\begin{document}
Este trabajo presenta un marco para el análisis de la calidad de pistas en sistemas interactivos de deducción basados en lenguaje natural. Para ello, se propone \textbf{QUILL} (\textit{Quality Index for Ludic Learning Hints}), una función compuesta que modela distintas dimensiones semánticas y cognitivas de las pistas, como su capacidad discriminativa, ambigüedad o ganancia de información.

Como entorno experimental, se ha desarrollado \textit{Spygame}, una plataforma interactiva que permite generar pistas automáticamente a partir de contenido estructurado (e.g., Wikipedia) y registrar la interacción de los usuarios. El sistema ha sido evaluado mediante más de 3.400 interacciones reales en dos escenarios diferenciados: un entorno de evocación libre (\textit{mundo abierto}) y un entorno de reconocimiento con espacio de búsqueda acotado (\textit{mundo cerrado}).

A partir de los datos recogidos, se analiza la relación entre las características de las pistas y su efectividad, permitiendo ajustar la función propuesta en función del comportamiento real de los usuarios. Los resultados muestran que la calidad de una pista depende principalmente de su capacidad para reducir el espacio de hipótesis de forma precisa, más que de su complejidad lingüística.
\end{document}
